% GENERATED FILE. Edit canonical lessons, then run npm run generate:books.
% canonical-source-hash: fnv1a64:40defa62728bf7e9
% canonical-lessons: GU-C151-notis, GU-C151-khullun, GU-C151-phakt, GU-C151-jaher, GU-C151-dukan-notis

\chapter{Reading a shop notice}
\label{ch:gu-reading-a-shop-notice}
\hlchaptermodality{151}

\begin{chapteropening}
\textbf{By the end of this chapter:} \emph{I can read a shop's notice of opening days and times, with its advert, and decide when I can go.}

Read a 45-word shop notice and decide when to go for rice, for milk, and on the public holiday.
\end{chapteropening}

\section[\texorpdfstring{\gu{નોટિસ} (noṭis)}{નોટિસ (noṭis)}]{\gu{નોટિસ} (noṭis) — a posted notice}

\label{lesson:GU-C151-notis}



\begin{quote}
Before the new one: say \emph{perhaps she'll come}, then \emph{I'm going to go to the market}, as a plan.
\end{quote}

\subsection*{You'll want to know: \gu{નોટિસ}}
\textbf{\gu{નોટિસ}} (\emph{noṭis}) — \textquotedblleft{}a posted notice, a notice\textquotedblright{}.

A notice is a sheet put up for everyone: on a shop door, at a station, at work. Gujarati took the English word whole.

\begin{quote}
\textbf{\gu{નોટિસ} \gu{ક્યાં} \gu{છે}?} — \emph{noṭis kyā̃ che?} — Where is the notice? \textbf{\gu{મેં} \gu{નોટિસ} \gu{વાંચી}.} — \emph{mẽ noṭis vā̃cī.} — I read the notice.
\end{quote}

\textbf{\gu{વાંચી}} ends in \textbf{-\gu{ી}} because \textbf{\gu{નોટિસ}} is feminine.

\begin{cousinweb}
\textbf{\gu{નોટિસ}} (\emph{noṭis}) is a loan from English \emph{notice}. English took it from French, and French from Latin \emph{notitia}, \textquotedblleft{}a being known\textquotedblright{}, built on \emph{nōscere}, \textquotedblleft{}to come to know\textquotedblright{}.
\end{cousinweb}

\begin{grammarlens}[title={What you've built}]
One more word for reading the notices on a street.
\end{grammarlens}

\subsection*{Your turn}
Say these aloud:

\begin{itemize}
\raggedright
  \item \emph{noṭis}
  \item \textbf{\gu{મેં} \gu{નોટિસ} \gu{વાંચી}} — \emph{mẽ noṭis vā̃cī}
  \item \emph{noṭis kyā̃ che?}
\end{itemize}

\begin{tcolorbox}[breakable,colback=teal!4,colframe=teal!35!black,title={Before you move on}]
What does \textbf{\gu{નોટિસ}} mean? (\emph{A posted notice}.) Why \textbf{\gu{વાંચી}} after it? (It is feminine.)
\end{tcolorbox}
\section[\texorpdfstring{\gu{ખુલ્લું} (khullũ)}{ખુલ્લું (khullũ)}]{\gu{ખુલ્લું} (khullũ) — not closed}

\label{lesson:GU-C151-khullun}



\begin{quote}
Before the new one: say \emph{a posted notice}, then \emph{I read the notice}.
\end{quote}

\subsection*{You'll want to know: \gu{ખુલ્લું}}
\textbf{\gu{ખુલ્લું}} (\emph{khullũ}) — \textquotedblleft{}not closed, open\textquotedblright{}.

It describes a door, a window or a shop that stands open, and it bends like \emph{my}:

\begin{quote}
\textbf{\gu{દુકાન} \gu{ખુલ્લી} \gu{છે}.} — \emph{dukān khullī che.} — The shop is open. \textbf{\gu{દરવાજો} \gu{ખુલ્લો} \gu{છે}.} — \emph{darvājo khullo che.} — The gate is open. \textbf{\gu{બારી} \gu{ખુલ્લી} \gu{હતી}.} — \emph{bārī khullī hatī.} — The window was open.
\end{quote}

\textbf{\gu{ખોલવું}} is the act, \emph{to open}; \textbf{\gu{ખુલ્લું}} is the state it leaves behind.

\begin{cousinweb}
\textbf{\gu{ખુલ્લું}} (\emph{khullũ}) is inherited. CDIAL sets \textbf{\gu{ખોલવું}} \emph{kholvũ}, \textquotedblleft{}to open\textquotedblright{}, and \emph{khulvũ}, \textquotedblleft{}to be open\textquotedblright{}, under a reconstructed Indo-Aryan \emph{khōll-} and \emph{khull-} (CDIAL 3945); \textbf{\gu{ખુલ્લું}} is the adjective of that second verb.
\end{cousinweb}

\begin{grammarlens}[title={What you've built}]
One more word for reading the notices on a street.
\end{grammarlens}

\subsection*{Your turn}
Say these aloud:

\begin{itemize}
\raggedright
  \item \emph{khullũ}
  \item \textbf{\gu{દુકાન} \gu{ખુલ્લી} \gu{છે}} — \emph{dukān khullī che}
  \item \textbf{\gu{દુકાન} \gu{ખુલ્લી} \gu{નથી}} — \emph{dukān khullī nathī}, the shop is not open
\end{itemize}

\begin{tcolorbox}[breakable,colback=teal!4,colframe=teal!35!black,title={Before you move on}]
What does \textbf{\gu{ખુલ્લું}} mean? (\emph{Not closed, open}.) Why \textbf{\gu{ખુલ્લી}} with \textbf{\gu{દુકાન}}? (\textbf{\gu{દુકાન}} is feminine.)
\end{tcolorbox}
\section[\texorpdfstring{\gu{ફક્ત} (phakt)}{ફક્ત (phakt)}]{\gu{ફક્ત} (phakt) — only}

\label{lesson:GU-C151-phakt}



\begin{quote}
Before the new one: say \emph{the shop is open}, then \emph{the window was open}.
\end{quote}

\subsection*{You'll want to know: \gu{ફક્ત}}
\textbf{\gu{ફક્ત}} (\emph{phakt}) — \textquotedblleft{}only, no more than\textquotedblright{}.

Notices are full of limits, and \textbf{\gu{ફક્ત}} sets them:

\begin{quote}
\textbf{\gu{દુકાન} \gu{ફક્ત} \gu{સવારે} \gu{ખુલ્લી} \gu{છે}.} — \emph{dukān phakt savāre khullī che.} — The shop is open only in the morning. \textbf{\gu{ફક્ત} \gu{પાંચ} \gu{રૂપિયા}.} — \emph{phakt pā̃c rūpiyā.} — Only five rupees.
\end{quote}

It stands right before the word it limits.

\begin{cousinweb}
\textbf{\gu{ફક્ત}} (\emph{phakt}) is usually traced, through Persian, to Arabic \emph{faqaṭ}, \textquotedblleft{}only, and that is enough\textquotedblright{}. Lane treats \emph{faqaṭ} under the root \emph{q-ṭ}, whose \emph{qaṭ} means \textquotedblleft{}enough, sufficient\textquotedblright{} (Lane).
\end{cousinweb}

\begin{grammarlens}[title={What you've built}]
One more word for reading the notices on a street.
\end{grammarlens}

\subsection*{Your turn}
Say these aloud:

\begin{itemize}
\raggedright
  \item \emph{phakt}
  \item \emph{dukān phakt savāre khullī che}
  \item \textbf{\gu{ફક્ત} \gu{ચા}} — \emph{phakt cā}, only tea
\end{itemize}

\begin{tcolorbox}[breakable,colback=teal!4,colframe=teal!35!black,title={Before you move on}]
What does \textbf{\gu{ફક્ત}} mean? (\emph{Only}.) Where does it stand? (Right before the word it limits.)
\end{tcolorbox}
\section[\texorpdfstring{\gu{જાહેર} (jāher)}{જાહેર (jāher)}]{\gu{જાહેર} (jāher) — public}

\label{lesson:GU-C151-jaher}



\begin{quote}
Before the new one: say \emph{only in the morning}, then \emph{only five rupees}.
\end{quote}

\subsection*{You'll want to know: \gu{જાહેર}}
\textbf{\gu{જાહેર}} (\emph{jāher}) — \textquotedblleft{}public, open to all\textquotedblright{}.

It marks what belongs to everyone, and on a notice it most often comes before \textbf{\gu{રજા}}, \emph{a holiday}:

\begin{quote}
\textbf{\gu{આવતીકાલે} \gu{જાહેર} \gu{રજા} \gu{છે}.} — \emph{āvtīkāle jāher rajā che.} — Tomorrow is a public holiday.
\end{quote}

A \textbf{\gu{જાહેર} \gu{રજા}} closes offices, banks and many shops for everybody at once.

\begin{cousinweb}
\textbf{\gu{જાહેર}} (\emph{jāher}) is usually traced, through Persian, to Arabic \emph{ẓāhir}, \textquotedblleft{}apparent, manifest\textquotedblright{}, from the root \emph{ẓ-h-r}, \textquotedblleft{}to come out into the open, to become evident\textquotedblright{} (Lane). What is out in the open is public.
\end{cousinweb}

\begin{grammarlens}[title={What you've built}]
One more word for reading the notices on a street.
\end{grammarlens}

\subsection*{Your turn}
Say these aloud:

\begin{itemize}
\raggedright
  \item \emph{jāher}
  \item \emph{āvtīkāle jāher rajā che}
  \item \textbf{\gu{શું} \gu{આજે} \gu{જાહેર} \gu{રજા} \gu{છે}?} — is today a public holiday?
\end{itemize}

\begin{tcolorbox}[breakable,colback=teal!4,colframe=teal!35!black,title={Before you move on}]
What does \textbf{\gu{જાહેર}} mean? (\emph{Public, open to all}.) Which word does it stand before on a notice? (\textbf{\gu{રજા}}.)
\end{tcolorbox}
\section[(reading)]{(reading) — a notice on a shop door}

\label{lesson:GU-C151-dukan-notis}



\begin{quote}
Say \emph{tomorrow is a public holiday}, then \emph{only in the morning}.

Nobody is talking to anybody on this page. A notice tells you what you may do, and when.
\end{quote}

\begin{tcolorbox}[breakable,skin=enhanced,colback=orange!6,colframe=orange!45!black,boxrule=0.5pt,arc=1mm,left=6pt,right=6pt,top=4pt,bottom=4pt,fonttitle=\bfseries,title={Read this}]
\begin{quote}
\textbf{\gu{નોટિસ}}
\end{quote}

\begin{quote}
\textbf{\gu{દુકાન} \gu{રોજ} \gu{સવારે} \gu{અને} \gu{સાંજે} \gu{ખુલ્લી} \gu{છે}.} \textbf{\gu{બપોરે} \gu{દુકાન} \gu{ખુલ્લી} \gu{નથી}.} \textbf{\gu{રવિવારે} \gu{રજા} \gu{છે}.} \textbf{\gu{આવતીકાલે}, \gu{સોમવારે}, \gu{જાહેર} \gu{રજા} \gu{છે}.} \textbf{\gu{મંગળવારે} \gu{દુકાન} \gu{ફક્ત} \gu{સવારે} \gu{ખુલ્લી} \gu{હશે}.} \textbf{\gu{બુધવારે} \gu{દુકાન} \gu{સવારે} \gu{અને} \gu{સાંજે} \gu{ફરી} \gu{ખુલ્લી} \gu{હશે}.} \textbf{\gu{આજે} \gu{કેરી} \gu{સસ્તી} \gu{છે}!} \textbf{\gu{સાંજે} \gu{દૂધ} \gu{અને} \gu{દહીં} \gu{પણ} \gu{છે}.} \textbf{\gu{આભાર}.}
\end{quote}

\emph{Read it:} the notice once, straight through, without translating

\emph{Read it:} it a second time, and find every word that says \emph{when}
\end{tcolorbox}

\subsection*{What to know first}
Almost every line opens on a \emph{when}: \textbf{\gu{રોજ}}, \textbf{\gu{બપોરે}}, then four days, each with the \textbf{-\gu{ે}} of \emph{on}. Two lines say \textbf{\gu{ખુલ્લી} \gu{હશે}}, \emph{will be open}, about days still to come. The line with \textbf{!} is an advert, set into the notice to bring you in.

\subsection*{Your turn}
\begin{itemize}
\raggedright
  \item \emph{Answer:} can you shop here tomorrow, Monday? (No: \textbf{\gu{જાહેર} \gu{રજા}}.)
  \item \emph{Answer:} milk on Wednesday evening: yes or no? (Yes: open, with milk in the evening.)
  \item \emph{Answer:} what is cheap today? (Mangoes: \textbf{\gu{કેરી}}.)
\end{itemize}

\begin{tcolorbox}[breakable,colback=teal!4,colframe=teal!35!black,title={Before you move on}]
Which day is the shop open only in the morning? (\textbf{\gu{મંગળવારે}}, Tuesday.) Which two days is it closed? (\textbf{\gu{રવિવારે}} and \textbf{\gu{સોમવારે}}: Sunday, and the public holiday.)
\end{tcolorbox}
